\subsection{ISA} \label{sec:design_isa}

The first step to creating an implementation of a \ac{cpu} is to define its behavior in the form of an \ac{isa}. For this work, several existing \acp{isa} have been explored, however none provide the desired combination of ease-of-use, visualisation quality and low cost. Therefore, in this chapter, a custom \ac{isa} is developed. The hardware built during this project then implements this architecture.

\subsubsection{Memory model}

As outlined in \autoref{sec:fundamentals_isa_memory_model}, there are 3 major memory models. For this project, a strict Harvard architecture is used since it allows a full separation of program instructions and general purpose data. This makes the timing required for memory accesses simpler, reducing the hardware needed to handle them. It also allows the architecture to use address buses of different size.

While most modern architectures have a word size of 64 bits or even more, this is impractical for the planned implementation due to the added number of components, wires and board space. On the other extreme, there are no architectures using exclusively 1 bit words. This would be too restrictive considering the program counter or the number of opcodes. However, there are architectures that use serial processing internally, operating on one bit at a time, such as the PDP8/S \cite{computermuseum_pdp8s}. There are several true 2-bit architectures such as the Intel 3000 \cite{wikichip_intel3000}. Similarly, there are historical architectures using words with a size of 4 bits each, such as the Intel 4004 and 4040 \cite{intel4004history}. However, a word size of 4 is impractical since even the most basic mathematical operations are likely to require operands greater than 15. Additionally, most memories and other similar components operate using 8 bit words. Since most students will also be familiar with the 8-bit byte, this is chosen as the word size for the architecture.

However, using 8 words for address data restricts the program counter and memory to 256 addressable locations. This means that any one program can only have at most 256 instructions, assuming each address holds one instruction. Likewise, there can only be 256 bytes of individually addressable \ac{ram}, assuming the memory is addressed on a byte-by-byte basis. While this can be useful to demonstrate the inner workings of a \ac{cpu}, most interesting real world problems require more storage than that. Since the \ac{isa} follows a strict Harvard architecture, the size of data and address words and busses can differ. In this case, using double the data word size as the address word size is a good compromise between available memory size and implementation cost. Using 16 bit addresses, a total of \SI{64}{\kibi\byte} can be addressed for both data \ac{ram} and program \ac{rom}. 



\subsubsection{Execution model}

The execution model of choice for this project is the register machine. Its working principles are explained in more detail in \autoref{sec:fundamentals_isa_execution_model}. Most modern machines are based on this architecture as well, due to better scalability and flexibility, as well as compilers being easier to optimize for such systems. Since the result of this work shall be used for demonstrating the working principles of digital computers, it makes sense to implement something close to real world systems. More specifically, the architecture of choice obeys the ``Load-store'' principle used for many modern RISC processors.

\subsubsection{Memory addressing modes}

% !TODO!